% IPIS Module 5 paper -- FLAT Editorial Manager variant (all sections inlined, no graphicspath).
% Generated from the split sources in paper5/. Upload this single file as the Manuscript;
% upload fig1_two_arm.png, fig2_deadtime.png, fig3_gamma_union.png and references.bib alongside.
% IPIS Module 5 (dynamic / horizon realization) paper for Computers & Chemical Engineering (Elsevier).
% Build: pdflatex main && bibtex main && pdflatex main x2
% WORKING TITLE -- provisional, pending Bien's ratification (see chat for alternatives).
\documentclass[review,12pt]{elsarticle}

\usepackage{amsmath,amssymb}
\usepackage{booktabs}
\usepackage{array}
\usepackage{graphicx}
\usepackage{lineno}
\usepackage{url}
\journal{Computers \& Chemical Engineering}

\begin{document}
\begin{frontmatter}

\title{Horizon-wide safety guarantees for closed-loop process operation via adaptive
conformal calibration}

\author{Bien Busico}
\address{Map\'ua Malayan College Mindanao, Davao City, Philippines}
\ead{bienbusico@gmail.com}

\begin{abstract}
% ---- begin abstract.tex ----
Process safety is judged over a campaign: what matters is the long-run frequency with which a
plant simultaneously meets its product specification and preserves its equipment to the next
maintenance window. We give a distribution-free guarantee on this horizon-wide joint event for a
closed loop with realistic dynamics and delayed measurements. A per-cycle version exists,
composing a soft-sensor budget, a prognostic budget, and the conditional-calibration budget of a
conformal real-time optimiser, but it assumes exchangeability and holds one cycle at a time.
Closed-loop feedback makes cycles dependent, laboratory labels arrive late, and the naive union
bound that would cover a horizon is vacuous. We realise the loop dynamically, as a validated
static gain wrapped in actuator, process, and analyser-deadtime transport lags, and lift the
per-cycle certificate with adaptive conformal inference. The conformal selection penalty is zero
under arbitrary label delay, so validity is preserved and dead time enters only the coverage
rate. On a debutanizer twin the health-constrained loop holds the joint event at its certified
floor of $0.75$ with a $95\%$ Wilson lower bound of $0.992$, whereas a health-blind optimiser
stays in specification on essentially every cycle yet over-refluxes until projected life
collapses, dropping the event to $0.008$. Coverage is invariant to label delays of zero to eight
cycles, and adaptive calibration keeps a finite interval at target coverage where the union
bound is vacuous beyond fifty cycles. The per-cycle certificate thus becomes a runtime,
horizon-wide safety monitor.
% ---- end abstract.tex ----
\end{abstract}

\begin{keyword}
adaptive conformal inference \sep horizon coverage \sep real-time optimisation \sep
predictive maintenance \sep distillation \sep closed-loop safety
\end{keyword}

\end{frontmatter}
\linenumbers

% ---- begin 01_intro.tex ----
\section{Introduction}
\label{sec:intro}

Closed-loop process operation is judged over a campaign, not a single cycle. What matters
operationally is the long-run frequency with which the plant simultaneously meets its product
specification and preserves its equipment to the next planned maintenance window. This paper
gives a distribution-free guarantee on exactly that horizon-wide joint event for a closed loop
with realistic dynamics and delayed measurements, and shows that it holds where the obvious
construction does not.

A per-cycle version of the guarantee already exists. The composed coverage certificate of
\citet{busico_m4} bounds the one-cycle joint event, the product staying in specification and
the equipment surviving, under exchangeable calibration residuals, by summing a soft-sensor
budget \citep{busico_m1}, a prognostic budget \citep{busico_m2}, and the conditional-calibration
budget of a conformal real-time optimiser \citep{busico_m3}. Two things break when the loop is
actually run. Feedback makes successive cycles dependent, so exchangeability fails; and
laboratory or analyser labels return several cycles late, so the residuals that calibrate the
bound are stale. The natural repair, a Bonferroni union bound that charges each cycle a fraction
of the error budget to cover $K$ cycles jointly, becomes vacuous as $K$ grows. A horizon
guarantee therefore needs a different mechanism, not a tighter bookkeeping of the per-cycle one.

We supply that mechanism by lifting the per-cycle certificate with adaptive conformal inference
\citep{gibbs2021}, and we realise the loop dynamically so the claim is tested against transport
lags and analyser dead time rather than asserted on a quasi-static surrogate. The contributions
are four.
\begin{enumerate}
\item \textbf{R1, a deployable horizon-wide safety certificate.} Adaptive calibration carries the
  per-cycle floor to the campaign: the conformal health-constrained loop holds the joint safety
  event at its certified floor over the horizon, with a $95\%$ Wilson lower bound of $0.992$
  against a floor of $0.75$.
\item \textbf{R2, the certificate's operational value, quantified.} A health-blind optimiser
  meets the quality specification on essentially every cycle yet over-refluxes until projected
  remaining useful life collapses, dropping the joint event to $0.008$, while the constrained
  loop holds it. Meeting a specification is not the same as operating safely, and the health
  channel with its budget is what closes the gap.
\item \textbf{R3, validity that is invariant to measurement dead time.} The conformal selection
  penalty stays zero under arbitrary label delay, so dead time moves into the coverage rate
  rather than breaking validity; the joint coverage stays at $1.000$ across delays of zero to
  eight cycles.
\item \textbf{R4, evidence that adaptivity is necessary.} Adaptive calibration keeps a finite
  interval at target coverage across horizons at which the naive union bound is already vacuous.
\end{enumerate}

This paper completes the IPIS series, which develops a physics-anchored soft sensor
\citep{busico_m1}, a similarity-calibrated prognostic certificate \citep{busico_m2}, a
conditionally calibrated conformal real-time optimiser \citep{busico_m3}, and the per-cycle
composed certificate \citep{busico_m4} that the present work lifts to a horizon.
Section~\ref{sec:framework} develops the framework: the per-cycle certificate, its dynamic
realisation, the delayed-label read-off, the dead-time argument, and the horizon guarantee. The
case study fixes the debutanizer twin and the two control arms, and Section~\ref{sec:results}
reports the campaign behaviour against the three results above. The discussion places the
guarantee as a runtime safety monitor and the bridge to plantwide risk allocation.
% ---- end 01_intro.tex ----
% ---- begin 02_related_work.tex ----
\section{Related work}
\label{sec:related}

Conformal prediction supplies finite-sample, distribution-free coverage under exchangeability
\citep{vovk2005, shafer2008}. Split and inductive variants make it cheap enough for online
regression \citep{lei2018, angelopoulos2023}, and conformalized quantile regression adapts the
interval width to local difficulty \citep{romano2019}. These deliver per-instance marginal
coverage, which is exactly the ingredient the per-cycle composed certificate of
\citet{busico_m4} sums across its soft-sensor and prognostic channels.

Closed-loop data are not exchangeable, because feedback couples successive cycles. Two lines of
work extend conformal validity past exchangeability. Reweighting recovers coverage under a
covariate shift that is known or estimable \citep{tibshirani2019}, and the
beyond-exchangeability framework bounds the coverage gap by a total-variation-type term for
general dependent data \citep{barber2023}. These results justify that a conformal interval need
not collapse when a control loop introduces dependence, but they do not on their own deliver a
long-run guarantee over a horizon under a feedback policy.

Adaptive conformal inference closes that loop directly. It treats the miscoverage level as a
state updated online from realised errors and guarantees a long-run coverage frequency however
the distribution drifts \citep{gibbs2021}. This long-run, shift-robust property is the engine of
the horizon guarantee developed here, and, as Section~\ref{sec:o1} shows, it tolerates the
delayed feedback that analyser dead time imposes.

On the control side, real-time optimisation drives a plant toward its economic optimum despite
plant-model mismatch, classically through modifier adaptation
\citep{marchetti2009, chachuat2009}, while self-optimizing and plantwide control frame which
variables to hold so that economic loss stays bounded \citep{skogestad2000}. The IPIS optimiser
\citep{busico_m3} adds a conformal back-off so that the economic push respects a chance
constraint; the present work asks what that back-off guarantees over a whole campaign rather
than at a single cycle.

This paper is the realisation layer of the IPIS series, which develops a physics-anchored soft
sensor \citep{busico_m1}, a similarity-calibrated prognostic certificate \citep{busico_m2}, the
conformal real-time optimiser \citep{busico_m3}, and the per-cycle composed certificate
\citep{busico_m4}. To our knowledge none of this prior work delivers a horizon-wide guarantee on
a joint quality-and-equipment-health event in a dynamic closed loop with delayed labels.
Establishing that guarantee, and showing that the naive union-bound alternative is vacuous at
realistic horizons, is the gap this paper fills.
% ---- end 02_related_work.tex ----
% ---- begin 03_framework.tex ----
\section{Framework}
\label{sec:framework}

\subsection{The per-cycle composed certificate}
\label{sec:percycle}

Each decision cycle $k$ produces a setpoint $u_k=(R_k,D_k)$, the reflux ratio and the
distillate rate. Under that setpoint the column realises a bottoms impurity $x_k$ and the
reflux pump carries a projected remaining useful life $\rho_k$. The operating requirement is
the \emph{joint} safety event
\begin{equation}
  S_k \;=\; \{\, x_k \le x_{\mathrm{spec}} \,\}\;\cap\;\{\, \rho_k \ge \rho_{\min} \,\},
  \label{eq:Sk}
\end{equation}
the product staying in specification \emph{and} the equipment surviving to its next planned
maintenance window. The composed coverage certificate of \citet{busico_m4} guarantees this
event one cycle at a time: under exchangeable calibration residuals and a similarity budget
$\varepsilon$,
\begin{equation}
  \Pr(S_k) \;\ge\; 1-(\alpha_1+\alpha_2)-\varepsilon,
  \label{eq:floor}
\end{equation}
where $\alpha_1$ is the soft-sensor miscoverage \citep{busico_m1}, $\alpha_2$ the prognostic
miscoverage \citep{busico_m2}, and $\varepsilon$ the conditional-calibration budget that the
conformal back-off of the real-time optimiser spends \citep{busico_m3}. We call the
right-hand side of \eqref{eq:floor} the certified floor. The bound is per cycle and rests on
exchangeability. The aim of this paper is to make it hold across an operating \emph{horizon}
once the loop is closed, so that successive cycles are dependent and the labels that calibrate
the bound arrive late.

\subsection{Dynamic realisation}
\label{sec:dynamic}

The certificate is established on a quasi-static twin, in which a setpoint maps instantly to
its steady-state outcome. A real column does not. We therefore retain the validated
steady-state response surface $g(\cdot)$ of the twin \citep{busico_m3}, fitted from the
process simulator, as a nonlinear static gain, and wrap it in linear transport dynamics. The
result is a Hammerstein-class plant with three transport elements, each a measured or
specified time constant rather than a free parameter:
\begin{enumerate}
\item an actuator first-order lag, so the realised setpoint $(\tilde R_k,\tilde D_k)$ relaxes
  toward the commanded $(R_k,D_k)$ with time constant $\tau_{\mathrm{act}}$;
\item a process first-order lag, so the impurity and the sensor temperature relax toward their
  steady-state targets $g(\tilde R_k,\tilde D_k)$ with constants $\tau_{\mathrm{proc}}$ and
  $\tau_{\mathrm{temp}}$;
\item an analyser dead time $\theta_a$ on the measured impurity, reflecting the
  gas-chromatograph cycle.
\end{enumerate}
By construction the settled steady state reproduces the twin exactly, so the calibration that
underlies \eqref{eq:floor} transfers without re-estimation. The loop runs on two timescales: a
fast integration step $\Delta t$ for the dynamics, and a slower decision interval
$T_{\mathrm{rto}}=n_{\mathrm{sub}}\,\Delta t$ at which a new setpoint is computed and then held.
Holding $u_k$ across the interval, rather than applying it instantaneously, is what makes the
similarity coordinate evolve along the realised path and what exposes the certificate to the
plant dynamics it must now survive.

\subsection{Delayed labels and the adaptive read-off}
\label{sec:adaptive}

A laboratory or analyser measurement of cycle $k$ does not return until $D_a$ decision cycles
later, where $D_a=\lceil \theta_{\mathrm{lab}}/T_{\mathrm{rto}}\rceil$. The horizon read-off
therefore replaces the per-cycle conformal quantile with an adaptive one
\citep{gibbs2021}. Writing $s_j$ for the soft-sensor nonconformity score of cycle $j$ and
$\mathrm{err}_j=\mathbf{1}[\,s_j>\hat q_j\,]$ for its miscoverage indicator, the adaptive level
follows
\begin{equation}
  \alpha_{t+1} \;=\; \alpha_t + \gamma\,\big(\alpha_1 - \mathrm{err}_{t}\big),
  \label{eq:aci}
\end{equation}
and the interval that the optimiser backs off against at cycle $k$ is
$\mathcal{C}_k=\hat y_k \pm \hat q\big(1-\alpha^{(k)}\big)$. Because the label for cycle $j$
is only consumed at cycle $j+D_a$, the level $\alpha^{(k)}$ that forms $\mathcal{C}_k$ depends
only on scores
\begin{equation}
  \{\, s_j : j \le k-1-D_a \,\},
  \label{eq:usable}
\end{equation}
that is, on residuals that are stale by $D_a$ cycles relative to the no-delay case. This is the
standard delayed-feedback regime for adaptive conformal inference, and it is the only channel
through which dead time enters the guarantee.

\subsection{Validity under dead time}
\label{sec:o1}

A closed loop raises a selection concern: if the optimiser were to choose $u_k$ using the same
sample later used to check coverage at cycle $k$, the marginal guarantee would no longer hold,
and \citet{busico_m4} carry an explicit selection penalty $\Delta_{\mathrm{sel},k}$ that is
non-zero only in that configuration. Dead time does not create such co-selection. By
\eqref{eq:usable} the interval $\mathcal{C}_k$ is a function of scores no newer than
$k-1-D_a$, and for every $D_a\ge 0$ this set excludes $s_k$ itself, since $k>k-1-D_a$. The
cycle-$k$ interval is thus formed without reference to the cycle-$k$ outcome it is scored
against, so the per-cycle result $\Delta_{\mathrm{sel},k}=0$ of \citet{busico_m4} is preserved
under arbitrary label delay. What dead time does change is that $\alpha^{(k)}$ is built from
older residuals; its effect is a transient in the coverage \emph{rate}, of order $D_a/K$ over a
$K$-cycle horizon, which washes out as the horizon grows, and not a loss of validity. The
detailed timing argument is given in Appendix~\ref{app:o1}.

\subsection{The horizon guarantee}
\label{sec:horizon}

Equations \eqref{eq:aci}--\eqref{eq:usable} deliver a long-run coverage property for the
soft-sensor interval that is distribution-free and holds under the dependence that feedback
induces \citep{gibbs2021}: the realised miscoverage frequency converges to $\alpha_1$ up to an
$O(1/(\gamma K))$ term, with an additional $O(D_a/K)$ boundary term under delay. Composing this
long-run quality guarantee with the prognostic and budget terms of \eqref{eq:floor} carries the
per-cycle floor to the horizon: the campaign frequency of the joint event $S_k$ meets the
certified floor of \eqref{eq:floor}. The naive alternative, a Bonferroni union bound that
charges each cycle the level $\alpha_1/K$ to cover $K$ cycles jointly, requires the
$1-\alpha_1/K$ residual quantile per cycle; as $K$ grows this level exceeds what any finite
calibration set can resolve and the interval becomes vacuous, so static union bounds cannot buy
horizon coverage. Section~\ref{sec:results} reports both behaviours.
% ---- end 03_framework.tex ----
% ---- begin 04_case_study.tex ----
\section{Case study}
\label{sec:case}

\subsection{Process twin}
The benchmark is a debutanizer column built in DWSIM \citep{dwsim} under the Peng-Robinson
equation of state \citep{peng1976}. From this twin we take the validated steady-state response
surface relating reflux ratio and distillate rate to the bottoms butane impurity
\citep{busico_m3} and use it as the nonlinear static gain of the dynamic plant; column
diagnostics follow standard correlations \citep{perrys}. Because the gain is the same surface
the per-cycle certificate was calibrated against, the settled steady state reproduces the twin
and the calibration transfers without re-estimation.

\subsection{Dynamic plant and health channel}
The static gain is wrapped in the transport dynamics of Section~\ref{sec:dynamic}. Representative
time constants are an actuator lag of about one minute on the manipulated setpoints, a
composition lag of about thirty minutes and a sensor-temperature lag of about fifteen minutes on
the outputs, and an analyser dead time of about five minutes on the measured impurity; the loop
is integrated at a one-minute step and a new setpoint is computed and held every sixty steps,
that is, every hour. The reflux pump degrades under an affinity cube law in the realised reflux,
and its projected remaining useful life is the equipment-health channel the certificate must
protect. These time constants are illustrative values, to be pinned per column from step-test
data; the coverage claim does not depend on them.

\subsection{Control arms}
Two arms share everything except a single switch. The health-blind arm lets the optimiser pull
reflux for economic benefit with no health constraint; the health-constrained arm caps that pull
with the conformal $\psi$-budget. The cap, not the conformal level, is what distinguishes the
arms, and both are scored against the same certified floor, so the contrast isolates the value of
enforcing the health budget.

\subsection{Configuration and reproducibility}
The frozen configuration is a bottoms-impurity specification of $0.02$, a minimum projected life
of $200$~h, soft-sensor and prognostic miscoverage $\alpha_1=\alpha_2=0.10$, and a similarity
budget $\varepsilon=0.05$, which set a certified floor of $0.75$. The two-arm contrast runs four
seeds over $120$ cycles; the delay and adaptivity sweeps run $400$ cycles, with a label delay of
two decision cycles unless that delay is itself swept. Every figure and number regenerates from a
single fixed seed through one script, and the headline numbers are frozen as committed evidence.
% ---- end 04_case_study.tex ----
% ---- begin 05_results.tex ----
\section{Results}
\label{sec:results}

We report a closed-loop campaign under the case-study configuration: bottoms-impurity
specification $x_{\mathrm{spec}}=0.02$, minimum projected life $\rho_{\min}=200$~h,
soft-sensor and prognostic miscoverage $\alpha_1=\alpha_2=0.10$, and similarity budget
$\varepsilon=0.05$, which set a certified floor of $0.75$. Every number below regenerates from
a single fixed seed and is frozen as committed evidence.

\subsection{The certificate's operational value}
\label{sec:res-twoarm}

Figure~\ref{fig:two_arm} contrasts the two control arms over the campaign, both scored against
the same floor. The health-blind optimiser keeps the product in specification on essentially
every cycle, with a quality coverage of $1.000$, but its unconstrained reflux pull drives the
pump hard: the projected remaining useful life clears $\rho_{\min}$ on only $0.008$ of cycles,
so the joint safety event collapses to $S_k=0.008$ (95\% Wilson interval $[0.003, 0.021]$) and
fails the floor. The conformal health-constrained arm holds $S_k=1.000$ with a Wilson lower
bound of $0.992$, meeting the floor, and its $\psi$-budget binds on every cycle ($1.00$) where
the blind arm never binds ($0.00$). Meeting a quality specification is therefore not the same
as operating safely: a loop can be in specification yet structurally unsafe over the horizon,
and the health channel together with the budget that enforces it is what carries a per-cycle
guarantee into one that survives the campaign.

\begin{figure}[t]
  \centering
  \includegraphics[width=0.82\linewidth]{fig1_two_arm.png}
  \caption{Two-arm contrast over an operating campaign, both arms scored against the same
  certified floor. The health-blind loop meets the quality spec but its over-reflux drives
  projected remaining useful life below the floor, so the joint safety event $S_k$ collapses;
  the conformal health-constrained loop holds $S_k$ at the floor.}
  \label{fig:two_arm}
\end{figure}

\subsection{Validity under measurement delay}
\label{sec:res-deadtime}

Figure~\ref{fig:deadtime} sweeps the label delay $D_a$. The joint coverage holds $S_k=1.000$
at every delay tested, $D_a\in\{0,1,2,4,8\}$ cycles, and the soft-sensor interval coverage
stays on its $0.90$ target throughout, moving only from $0.897$ at $D_a=0$ to $0.895$ at
$D_a=8$. Delay leaves validity intact and shows up only as a small movement in the coverage
rate, consistent with the $\Delta_{\mathrm{sel},k}=0$ argument of Section~\ref{sec:o1}: the
cost of stale feedback is paid in the rate, not in validity.

\begin{figure}[t]
  \centering
  \includegraphics[width=0.82\linewidth]{fig2_deadtime.png}
  \caption{Horizon coverage versus label delay $D_a$. The joint-safety coverage $S_k$ holds the
  certified floor at every delay, and the soft-sensor interval coverage stays on target;
  validity is invariant to analyser and laboratory dead time.}
  \label{fig:deadtime}
\end{figure}

\subsection{Adaptive versus static calibration}
\label{sec:res-gamma}

Figure~\ref{fig:gamma_union}(a) varies the adaptive learning rate $\gamma$: interval coverage
stays near the $0.90$ target while the half-width tightens as $\gamma$ rises, and the
one-standard-error rule selects $\gamma=0.128$ at a mean half-width of $0.0011$.
Figure~\ref{fig:gamma_union}(b) sets this against the naive alternative. A $K$-cycle Bonferroni
interval widens from a half-width of $0.0009$ at $K=1$ to $0.0013$ at $K=10$ and is already
vacuous at $K\ge 50$, where no finite calibration quantile resolves the per-cycle level
$\alpha_1/K$; adaptive conformal inference instead holds a finite half-width at target coverage
across the same horizons \citep{gibbs2021}. Horizon coverage thus requires adaptivity, and the
static union bound cannot deliver it, which is the mechanism behind the horizon guarantee of
Section~\ref{sec:horizon}.

\begin{figure}[t]
  \centering
  \includegraphics[width=0.98\linewidth]{fig3_gamma_union.png}
  \caption{(a) Coverage versus adaptivity across the adaptive-conformal learning rate. (b) The
  naive $K$-cycle Bonferroni half-width becomes vacuous at realistic horizons, whereas adaptive
  conformal inference maintains a finite half-width at target coverage.}
  \label{fig:gamma_union}
\end{figure}
% ---- end 05_results.tex ----
% ---- begin 06_discussion.tex ----
\section{Discussion}
\label{sec:discussion}

Under adaptive calibration the per-cycle certificate of \citet{busico_m4} becomes a guarantee
that holds over the operating horizon rather than at an instant. In practice this turns the
certificate into a runtime monitor: at any point in a campaign the loop can state a
distribution-free lower bound on the frequency with which it has met the joint safety event,
and that bound does not weaken as the campaign lengthens (R1, Section~\ref{sec:horizon}).

The two-arm contrast in Figure~\ref{fig:two_arm} makes the operational content concrete, and it
is the result most likely to matter on a plant. A health-blind optimiser is in specification on
essentially every cycle and would pass any quality-only audit, yet it over-refluxes until
projected life collapses and the joint safety event falls to $0.008$, while the conformal
health-constrained loop holds the event at its floor. A quality specification and a safety
guarantee are different objects, and the binding $\psi$-budget is the mechanism that converts a
per-cycle bound into one sustained across the horizon (R2). The gap here is not subtle: it is
$0.008$ against $1.000$ on the same floor.

Delayed feedback is the feature of a real analyser loop most often assumed fatal to a conformal
guarantee, and it is not. Because the interval at a cycle never references that cycle's own
outcome, the selection penalty stays zero under any label delay (Section~\ref{sec:o1}), so
coverage remains valid; the only cost is a transient in the coverage rate, of order $D_a/K$,
which Figure~\ref{fig:deadtime} shows is small over the horizons tested (R3). It is worth
separating the property that survives dead time, validity, from the one that merely degrades
gracefully with it, the rate.

The contrast with the union bound in Figure~\ref{fig:gamma_union} is the methodological point. A
static Bonferroni allocation cannot buy horizon coverage without charging each cycle a vanishing
error level, at which point no finite calibration set resolves the required quantile and the
interval is vacuous; adaptive inference sidesteps this by spending a fixed long-run budget
rather than a per-cycle one \citep{gibbs2021}. Adaptivity is therefore not a refinement but a
requirement at realistic horizons (R4).

Several limitations bound these claims. The plant is a reduced-order Hammerstein realisation, a
validated static gain wrapped in transport lags rather than a tray-resolved dynamic column, so
the transient shapes are representative rather than exact and the time constants are
illustrative values to be pinned per column from step-test data. The coverage claim is itself
fidelity-agnostic: validity follows from the calibration construction, not from the plant's
internal accuracy, so what a higher-fidelity dynamic simulation would sharpen is the realism of
the coverage-rate transient, not the validity. A confirmation on a DWSIM dynamic column is the
natural next step. Finally, the deadtime transient is small here because the loop is
well-behaved; a deliberate drift or disturbance stress test would make the rate cost visible and
is left to future work.

The same composition that carries one column's certificate across a horizon points to the larger
target: allocating the similarity budget $\varepsilon$ across several units, so that a plantwide
certificate is assembled from per-unit ones under a shared risk allocation. That modular
extension, from one column to a plant, is the highest-leverage direction this framework opens.
% ---- end 06_discussion.tex ----
% ---- begin 07_conclusion.tex ----
\section{Conclusion}
\label{sec:conclusion}

We have given a horizon-wide safety guarantee for closed-loop process operation, obtained by
lifting the per-cycle composed coverage certificate of \citet{busico_m4} with adaptive conformal
inference \citep{gibbs2021} and testing it on a dynamically realised debutanizer loop with
delayed measurements. The guarantee is distribution-free and holds over an operating campaign
rather than at a single cycle.

Four results support it. The conformal health-constrained loop holds the joint safety event at
its certified floor over the horizon, with a $95\%$ Wilson lower bound of $0.992$ against a floor
of $0.75$. A health-blind optimiser meets the quality specification on essentially every cycle
yet over-refluxes until projected life collapses, dropping the joint event to $0.008$, which
shows that meeting a specification and operating safely are distinct and that the binding health
budget is what closes the gap. Validity is invariant to measurement dead time, the selection
penalty staying zero under any label delay, so dead time enters the coverage rate rather than
breaking the guarantee. And adaptive calibration retains a finite interval at target coverage
across horizons at which the naive union bound is vacuous, so adaptivity is necessary rather than
optional.

Taken together these carry the IPIS stack from a certified collection of components to a
certified operating campaign: the per-cycle certificate becomes a runtime, horizon-wide monitor.
The natural next steps are to allocate the similarity budget across several units, assembling a
plantwide certificate from per-unit ones under a shared risk allocation, and to confirm the
guarantee on an operating LPG or petrochemical column with live distributed-control-system data.
% ---- end 07_conclusion.tex ----

\section*{Data and code availability}
The dynamic plant, the dynamic closed-loop orchestrator, the adaptive-conformal horizon
read-off, the two-arm and sweep experiments, and the single script that regenerates every
figure in this paper (\texttt{scripts/run\_horizon\_sweeps.py}) are available at
\url{github.com/beebzy-droid/IPIS}; all figures and the reported coverage numbers reproduce
from a single fixed seed, with the headline numbers frozen as committed evidence. The
process twin is built in DWSIM \citep{dwsim} under the Peng-Robinson equation of state
\citep{peng1976}; the steady-state response surface it corrects, the soft sensor, the
calibrated prognostics, and the per-cycle composed coverage certificate that this work lifts
to a horizon are developed in the companion papers \citep{busico_m3, busico_m1, busico_m2, busico_m4}.

\section*{Declaration of competing interest}
The author declares no competing financial interests or personal relationships that
could have influenced the work reported in this paper.

\appendix
% ---- begin 08_appendix_o1.tex ----
\section{Timing of the delayed-label read-off and the zero selection penalty}
\label{app:o1}

This appendix records the timing argument behind Section~\ref{sec:o1}: the conformal selection
penalty $\Delta_{\mathrm{sel},k}$ stays zero under arbitrary label delay, so dead time affects
the coverage rate and not the validity of the certificate.

\paragraph{Timing of one decision cycle.}
The setpoint $u_k=(R_k,D_k)$ is computed at the end of cycle $k-1$, held across the decision
interval, after which the realised outcome $x_k$ with measurement noise $\varepsilon_k$ follows.
The laboratory or analyser reports the quality of cycle $k$ only
$D_a=\lceil \theta_{\mathrm{lab}}/T_{\mathrm{rto}}\rceil$ decision cycles later, and that delayed
label drives the adaptive update of \eqref{eq:aci}. The adaptive level $\alpha^{(k)}$ that forms
the cycle-$k$ interval $\mathcal{C}_k$ reflects every label applied through the end of cycle
$k-1$; since the label for cycle $j$ is applied at cycle $j+D_a$, the scores that influence
$\alpha^{(k)}$ are exactly $\{\, s_j : j \le k-1-D_a \,\}$.

\paragraph{The selection penalty vanishes.}
$\mathcal{C}_k$ is a deterministic function of $\alpha^{(k)}$, which depends only on
$\{\, s_j : j \le k-1-D_a \,\}$. For every $D_a\ge 0$ this set excludes the cycle-$k$ score
$s_k$, and with it $x_k$ and $\varepsilon_k$, because
\begin{equation*}
  k \;>\; k-1-D_a \qquad \text{for all } D_a \ge 0 .
\end{equation*}
The cycle-$k$ interval is therefore formed without reference to the cycle-$k$ outcome it is
scored against. Combined with the causal-timing property of the per-cycle certificate, that
$u_k$ is fixed before $\varepsilon_k$ and is hence independent of it given the operating regime
\citep{busico_m4}, the quantity selected at cycle $k$ does not condition on the realisation
against which it is checked. Hence $\Delta_{\mathrm{sel},k}=0$, and the per-cycle composed
certificate applies cycle by cycle exactly as in the quasi-static case; dead time does not
re-open the co-selection window.

\paragraph{Where the cost goes.}
What dead time changes is that $\alpha^{(k)}$ is built from scores no newer than $k-1-D_a$, that
is, stale by $D_a$ cycles relative to the no-delay case. The long-run coverage guarantee of
adaptive conformal inference bounds the time-averaged miscoverage by an $O(1/(\gamma K))$ term
that is independent of the data-generating dependence \citep{gibbs2021}; delayed feedback
re-indexes the error stream by $D_a$, and the same averaging leaves the long-run coverage at its
target up to an additive $O(D_a/K)$ boundary term that vanishes as $K\to\infty$. Long-run
coverage therefore holds, while the finite-horizon deviation carries a transient that grows with
$D_a$. That transient is the object characterised by the dead-time sweep of
Section~\ref{sec:results}.
% ---- end 08_appendix_o1.tex ----

\bibliographystyle{elsarticle-num-names}
\bibliography{references}

\end{document}
